% =====================================================================
%  Thesis proposal, Appendix: Concentration along the time axis
%  Owns: apx:time
% =====================================================================

\section{Concentration along the time axis}
\label{apx:time}

This appendix supports \cref{sec:profile}. It is a sketch, and it says
where it stops. Two elementary bounds and the compactness they give are
proved; the application of Lions' lemma is a statement about a sequence of
densities and is exact; the identification of its alternatives with the
flight regimes, and the passage from the dissipation back to the state,
are stated as what remains to be shown, and gathered in
\cref{sec:time-open}.

%--------------------------------------------------------------------------
\subsection{Setting, and two elementary bounds}\label{sec:time-setting}
%--------------------------------------------------------------------------

Throughout, $\eps = \eps_{\mathrm{atm}}$. The corridor $\calX$ is a compact
set of states on which $V \ge V_{\min} > 0$,
$\abs{\gamma}\le\gamma_{\max}<\pi/2$ and $\abs{\phi}\le\phi_{\max}<\pi/2$.
A \emph{trajectory} is an absolutely continuous
$x_\eps\colon[0,T]\to\calX$ whose translational part
satisfies~\crefrange{eq:rdot}{eq:phidot} and~\crefrange{eq:Vdot}{eq:psidot}
for almost every~$t$, with the specific aerodynamic force
$(F_V,F_\gamma,F_\psi)/m$ selected measurably from the enclosure. Nothing is
assumed about how the force arises: the bank history, the attitude and the
selection from the interval field are all free. One assumption is made on
the enclosure.

\begin{assum}[Dissipative enclosure]\label{asm:dissipative}
There is a constant $E_{\max}$ such that every selection satisfies
$F_V \le 0$ and $(F_\gamma^2 + F_\psi^2)^{1/2} \le E_{\max}\abs{F_V}$.
\end{assum}

The first inequality says that the aerodynamic force never does positive
work on the vehicle, the second that lift is bounded by a multiple of drag.
For a lift-and-drag enclosure they read $[C_D]\subset(0,\infty)$ and
$[C_L]/[C_D]\subset[-E_{\max},E_{\max}]$. A propulsive mode adds a bounded
source to what follows and is not considered here.

\begin{lem}[Energy balance]\label{lem:energy}
Let
\[
  \mathcal{E} \;=\; \tfrac12 V^2 + \Phi_g(r,\phi)
                    - \tfrac12\,\omega_E^2 r^2\cos^2\!\phi
\]
be the specific mechanical energy in the rotating frame, with $\Phi_g$ the
potential of \cref{sec:gravity}. Along every trajectory
\begin{equation}\label{eq:dissipation}
  \dot{\mathcal{E}} \;=\; \frac{F_V}{m}\,V \;=:\; -P \;\le\; 0 ,
\end{equation}
so that $\int_0^T P\dd t = \mathcal{E}(x(0)) - \mathcal{E}(x(T))
\le \Delta\mathcal{E}$, where
$\Delta\mathcal{E} = \max_{\calX}\mathcal{E} - \min_{\calX}\mathcal{E}$.
\end{lem}

\begin{proof}
Differentiate along \crefrange{eq:rdot}{eq:phidot} and \cref{eq:Vdot}. The
terms in $g_D$ and $g_N$ cancel against
$\partial_r\Phi_g\,\dot r + \partial_\phi\Phi_g\,\dot\phi$, by the
definition of $g_D$ and $g_N$ as the gradient of $\Phi_g$, and the
centrifugal term of \cref{eq:Vdot} against the derivative of
$\tfrac12\omega_E^2r^2\cos^2\!\phi$. The Coriolis force and the two normal
components of the aerodynamic force are perpendicular to the velocity and
do not appear in \cref{eq:Vdot} at all. The identity is checked by computer
algebra with oblateness and rotation retained.
\end{proof}

\begin{lem}[Impulse bound]\label{lem:impulse}
Under \cref{asm:dissipative} the aerodynamic acceleration
$\vvec{a} = \vvec{F}_a/m$ satisfies $\abs{\vvec{a}} \le C_a P$ with
$C_a = \sqrt{1+E_{\max}^2}\,/\,V_{\min}$. Hence for every interval
$I\subseteq[0,T]$,
\[
  \int_I \abs{\vvec{a}}\dd t \;\le\; C_a\int_I P\dd t .
\]
\end{lem}

\begin{proof}
$\abs{\vvec{a}}^2 = (F_V^2 + F_\gamma^2 + F_\psi^2)/m^2
\le (1+E_{\max}^2)\,F_V^2/m^2$ and $\abs{F_V}/m = P/V \le P/V_{\min}$.
\end{proof}

The second lemma is what makes the dissipation the right quantity to follow.
Lift cannot be had without drag, so wherever no energy is lost the
aerodynamics does nothing, and the motion is that of the gravity field in
the rotating frame. Whatever a pass does to the state, the dissipation sees
it.

%--------------------------------------------------------------------------
\subsection{Compactness in slow time}\label{sec:time-slow}
%--------------------------------------------------------------------------

Let $\eps_n\to 0$ and let $x_n$ be a trajectory for each~$n$. Write
$\mathcal{E}_n(t) = \mathcal{E}(x_n(t))$; $\pi_n = P_n\dd t$, a positive
measure on $[0,T]$ of mass at most $\Delta\mathcal{E}$; and
$\vvec{y}_n(t)$, $\vvec{v}_n(t)\in\R^3$ for the position and the
planet-relative velocity in planet-fixed Cartesian components, in which the
chart singularities play no part.

\begin{prop}[Skeleton and impulses as a limit]\label{prop:slow}
Under \cref{asm:dissipative} there is a subsequence along which:
\begin{enumerate}[label=(\roman*),leftmargin=*,itemsep=2pt]
  \item $\mathcal{E}_n(t)\to\mathcal{E}_0(t)$ for every $t\in[0,T]$, with
    $\mathcal{E}_0$ non-increasing, and $\pi_n\rightharpoonup\pi_0$
    weakly-$*$, where $\pi_0$ is the measure with distribution function
    $\mathcal{E}_0(0)-\mathcal{E}_0$;
  \item $\pi_0 = p\dd t + \sum_j a_j\,\delta_{t_j} + \pi_c$, with
    $p\in L^1(0,T)$, at most countably many atoms $a_j>0$, and $\pi_c$
    singular and without atoms; and for every $\delta>0$
    \[
      \#\{\,j : a_j \ge \delta\,\} \;\le\; \Delta\mathcal{E}/\delta ;
    \]
  \item on every closed interval on which $\mathcal{E}_0$ is continuous,
    the convergence $\mathcal{E}_n\to\mathcal{E}_0$ is uniform;
  \item $\vvec{y}_n\to\vvec{y}_0$ uniformly, $\vvec{v}_n(t)\to\vvec{v}_0(t)$
    for every~$t$, $\dot{\vvec{y}}_0 = \vvec{v}_0$ almost everywhere, and
    for $0\le s<t\le T$
    \begin{equation}\label{eq:bv}
      \abs{\vvec{v}_0(t) - \vvec{v}_0(s)}
      \;\le\; C_g\,(t-s) + C_a\,\pi_0([s,t]) ,
    \end{equation}
    with $C_g$ a bound on the gravitational, Coriolis and centrifugal
    accelerations over the corridor.
\end{enumerate}
\end{prop}

\begin{proof}
(i) The functions $\mathcal{E}_n$ are non-increasing by \cref{lem:energy}
and uniformly bounded on the compact corridor, so Helly's selection theorem
gives a subsequence converging at every point; the measures $\pi_n$ are
their Stieltjes measures, and pointwise convergence of the distribution
functions at the continuity points of the limit, with bounded total mass,
is weak-$*$ convergence. (ii) is the Lebesgue decomposition of $\pi_0$, and
the count follows from $\sum_j a_j\le\pi_0([0,T])\le\Delta\mathcal{E}$.
(iii) is P\'olya's theorem: monotone functions converging pointwise to a
continuous limit on a compact interval converge uniformly. (iv) In
planet-fixed components Newton's law reads
$\dot{\vvec{v}}_n = \vvec{a}_n + \vvec{g}(\vvec{y}_n)
 - 2\,\vvec{\omega}_E\times\vvec{v}_n
 - \vvec{\omega}_E\times(\vvec{\omega}_E\times\vvec{y}_n)$,
so by \cref{lem:impulse}
\[
  \abs{\vvec{v}_n(t)-\vvec{v}_n(s)}
  \;\le\; C_g\,(t-s) + C_a\,\pi_n([s,t]) .
\]
The velocities therefore have total variation at most
$C_g T + C_a\,\Delta\mathcal{E}$, uniformly in~$n$, and Helly's theorem for
functions of bounded variation gives pointwise convergence along a further
subsequence. The positions are uniformly Lipschitz and converge uniformly
by Arzel\`a--Ascoli, and $\dot{\vvec{y}}_0=\vvec{v}_0$ by dominated
convergence. \Cref{eq:bv} follows on letting $n\to\infty$, since
$\limsup_n\pi_n(F)\le\pi_0(F)$ for closed~$F$.
\end{proof}

Three things can be read off \cref{eq:bv}. The limiting velocity has bounded
variation, and it can jump only at an atom of $\pi_0$, by at most
$C_a a_j$: an atom is an impulse, and its size bounds the impulse. On an
interval that $\pi_0$ does not charge, $\int\abs{\vvec{a}_n}\dd t\to 0$ and
the limit is an arc of the gravity field alone. And where $\pi_0$ has a
density~$p$ the aerodynamic forcing survives as an acceleration of
magnitude at most $C_a\,p$, which is the glide. This is the
skeleton-and-impulses picture as a statement about
limits, valid for every control and every selection, and it is as far as
compactness alone goes. It does not exclude the part $\pi_c$, and it says
nothing yet about what happens \emph{inside} an atom.

%--------------------------------------------------------------------------
\subsection{The fast-time axis and Lions' lemma}\label{sec:time-fast}
%--------------------------------------------------------------------------

To see inside an atom the time axis is dilated. Fix a scale $h_n\to 0$---in
the body, $h_n=\sqrt{\eps_n}$---suppose the total dissipation stays bounded
below, $\pi_n([0,T])\ge m_0>0$, and set
\[
  \nu_n(\tau) \;=\; \frac{h_n\,P_n(h_n\tau)}{\pi_n([0,T])}
  \qquad(\tau\in\R),
\]
extended by zero outside $[0,T/h_n]$. Then $\nu_n\ge 0$ and
$\int_\R\nu_n\dd\tau = 1$, and the concentration-compactness
lemma~\cite[Lemma~I.1]{lions1984concentration1} applies with no further
hypothesis. There is a subsequence along which one of the following holds.
\begin{description}[leftmargin=1.5em,itemsep=2pt]
  \item[Compactness.] There are $\tau_n\in\R$ such that for every
    $\eta>0$ there is $R<\infty$ with
    $\int_{\tau_n-R}^{\tau_n+R}\nu_n\dd\tau \ge 1-\eta$ for all~$n$.
  \item[Vanishing.] For every $R<\infty$,
    $\;\sup_{y\in\R}\int_{y-R}^{y+R}\nu_n\dd\tau \to 0$.
  \item[Dichotomy.] There is $\alpha\in(0,1)$ such that for every $\eta>0$
    there are nonnegative $\nu_n^1,\nu_n^2$ with
    $\norm{\nu_n - \nu_n^1 - \nu_n^2}_{L^1}\le\eta$,
    $\abs{\int\nu_n^1 - \alpha}\le\eta$,
    $\abs{\int\nu_n^2 - (1-\alpha)}\le\eta$, and
    $\operatorname{dist}(\operatorname{supp}\nu_n^1,
    \operatorname{supp}\nu_n^2)\to\infty$.
\end{description}

Each has a reading. In the first, all but a fraction $\eta$ of the
dissipation occurs within a time $R\,h_n$ of $t_n = h_n\tau_n$: a single
pass, and $\pi_0$ is one atom. In the second no window of width $R\,h_n$
holds a fixed share. This is the case whenever the drag power is bounded
uniformly in~$n$, since then a window holds at most $2RC\,h_n/m_0$; the
limit $\pi_0$ then has a bounded density and no atom. In the third the
dissipation splits into pieces that separate on the fast scale: two passes
a positive time apart, or a pass followed by a glide. Applying the lemma
again to each piece, as in the construction of a profile decomposition,
extracts the concentrated pieces one at a time; for a threshold~$\delta$
the process stops after at most $\Delta\mathcal{E}/\delta$ steps, by
\cref{prop:slow}(ii), and leaves a remainder in which no window holds more
than about~$\delta$. The remainder is small in the sense of the second
alternative, in the norm $\sup_y\int_{y-R}^{y+R}\abs{\,\cdot\,}\dd\tau$,
and not in $L^1$: on a glide it is most of the dissipation.

\begin{remark}[Which alternative occurs is decided by the corridor]
\label{rem:regimes}
Whether a family concentrates at all is a statement about how its corridor
scales with $\eps$, and belongs among the hypotheses. If the dynamic
pressure and load limits are held fixed as $\eps\to 0$, the drag power is
bounded, every sequence vanishes, and the limit has no atoms: the
trajectories are glide and phugoid, and the fast-scale structure is
\emph{oscillation}, for which the natural tool is averaging over the
invariant measures of the fast flow~\cite{artstein1999invariant}.
Concentration requires loads that grow as $\eps\to 0$, as in the classical
matched-asymptotic analysis of skip entry at a fixed entry
angle~\cite{shi1971skip}, where the layer has thickness $O(\eps)$ in
altitude. The physical value $\eps\approx 1.1\times10^{-3}$, with entry
angles of a few degrees, lies between the two: by \cref{eq:turning} a pass
that turns the flight path through $8^\circ$ at a vertical lift-to-drag
ratio of two dissipates about an eighth of the kinetic energy, and a
computed pass lasts about a fifth of the slow time.
The threshold in \cref{eq:kmax} is what makes the decomposition well posed
at fixed $\eps$ without deciding the regime. Passes that dissipate more
than $\delta$ are bubbles, and there are at most $\Delta\mathcal{E}/\delta$
of them whatever $\eps$ is; everything else is residual, and has to be
bounded as oscillation.
\end{remark}

\begin{remark}[The scale has to be chosen for each piece]
Compactness at the scale $h_n$ forces a single atom of $\pi_0$, but
vanishing at that scale does not exclude one: dissipation that concentrates
at a coarser scale $h_n'\gg h_n$, $h_n'\to 0$, vanishes at the scale $h_n$
and still produces an atom. Conversely a pass shorter than $h_n$ is compact
at that scale and its profile is a Dirac mass. The lemma classifies what
happens at one scale; the width of each piece has to be found from the
piece. That is the content of the limit-case form of the
principle~\cite{lions1985limit1}, in which compactness is recovered up to a
translation \emph{and} a dilation.
\end{remark}

%--------------------------------------------------------------------------
\subsection{Profiles and their orthogonality}\label{sec:time-profiles}
%--------------------------------------------------------------------------

For an atom $(t_j,a_j)$ of $\pi_0$ choose centers $t_n^j\to t_j$ and widths
$h_n^j\to 0$ such that the windows
$[t_n^j - R\,h_n^j,\, t_n^j + R\,h_n^j]$ capture $a_j$ as $R\to\infty$,
uniformly in~$n$, with $h_n^j$ the smallest scale at which this holds. The
rescaled densities $s\mapsto h_n^j\,P_n(t_n^j + h_n^j s)$ are then tight on
the line and converge weakly-$*$ along a subsequence to a measure $p^j$ of
mass $a_j$, the \emph{profile} of the pass. Two profiles $j\ne k$ are
\emph{orthogonal} in the sense of~\cite[(1.18)]{bahouri1999high} when
\[
  \Bigl|\log\frac{h_n^j}{h_n^k}\Bigr| \to \infty
  \qquad\text{or}\qquad
  h_n^j = h_n^k \ \text{ and } \
  \frac{\abs{t_n^j - t_n^k}}{h_n^j}\to\infty .
\]
Atoms at distinct times are orthogonal by the second alternative, with
nothing to prove. For orthogonal profiles the dissipation is additive,
\begin{equation}\label{eq:additivity}
  \pi_n([0,T])
  \;=\; \sum_{j=1}^{J} a_j \;+\; r_n^J([0,T]) \;+\; o(1)
  \qquad(n\to\infty),
\end{equation}
with $r_n^J$ the dissipation outside the $J$ windows. This is the analogue
of the energy identity~\cite[(1.25)]{bahouri1999high}, and it is immediate,
because the dissipation is a measure and the windows are eventually
disjoint. The analogue of the nonlinear
statement~\cite[(1.26)]{bahouri1999high}---that the solution is, up to a
small error, the superposition of the profiles' separate evolutions---is
not immediate, and here it is not a superposition at all. An ordinary
differential equation composes: the state leaving pass~$j$ is the datum of
the arc that leads to pass~$j+1$. What replaces (1.26) is a factorization
of the map from entry data to exit data through the individual pass maps
and the outer flow, with a remainder that is explicit and uniform over the
selections. That is objective~4.

%--------------------------------------------------------------------------
\subsection{Joint profiles where scales meet}\label{sec:time-joint}
%--------------------------------------------------------------------------

Within a single pass the attitude and elastic motions are excited at their
own scales, $h^{\mathrm{att}} = \eps_{\mathrm{att}}$ and
$h^{\mathrm{ela}} = \eps_{\mathrm{ela}}$ in units of the slow time. They
share the center of the pass, so the second alternative of (1.18) is not
available and orthogonality has to come from the first: the ratio of the
scales must diverge along the family. Where it does, the faster motion
averages out over the slower. Where two periods remain comparable---where
two bands cross inside the corridor---the corresponding profiles cannot be
separated, and the decomposition carries one joint profile on the coupled
subsystem. The check is pairwise and depends on the vehicle. For generic
parametrizations the first structural period lies between
$6\,\mathrm{ms}$ and $0.18\,\mathrm{s}$ and a pass lasts about
$150\,\mathrm{s}$; with the attitude period taken between $0.1$ and
$10\,\mathrm{s}$, the atmospheric scale is separated from the other two by
more than a decade, and the attitude and elastic scales meet only for a
very slender, flexible body with a fast attitude loop. These are orders of
magnitude, to be redone for any vehicle to which the result is applied.

%--------------------------------------------------------------------------
\subsection{What is not proved}\label{sec:time-open}
%--------------------------------------------------------------------------

\begin{enumerate}[leftmargin=*,itemsep=3pt]
  \item \emph{The singular continuous part.} \Cref{prop:slow} allows
    $\pi_c\ne 0$: dissipation carried by a set of times of measure zero
    without atoms. Countably many passes that accumulate at a finite time,
    with summable dissipation---a sequence of decaying skips that ends in a
    glide, as the bounces of a ball end in rolling---are atoms and are
    already covered. A genuinely diffuse singular part would have to be
    excluded by a lower bound on the separation of passes in terms of the
    energy each dissipates. This step is open.
  \item \emph{Vanishing is glide.} That the part of $\pi_0$ with a density
    corresponds to motion near the glide balance, and how near, is the
    estimate of \cref{apx:residual}. It is an averaging estimate for a
    weakly damped oscillator, not a boundary-layer estimate, and it is not
    yet uniform over the selections.
  \item \emph{Compactness of the profiles over the family.} The profile
    $p^j$ and the impulse set it generates depend on the selection.
    \Cref{lem:impulse} bounds the diameter of the impulse set by $C_a a_j$;
    a description of the set itself is the inner problem of the classical
    theory, to be solved over the interval field.
  \item \emph{Composition.} The factorization that replaces (1.26), with
    its remainder, at each of the three scales and under interval
    disturbances.
\end{enumerate}
